\chapter[Hľadanie politiky správania]{Hľadanie politiky správania a tréning agenta}

\section{Problém: hľadáme politiku správania}

Predchádzajúca kapitola pripravila rozhranie medzi hrou a agentom. Máme stav auta, virtuálnu geometriu trate, vektor pozorovania, akčný priestor a spôsob, ako akciu poslať späť do hry. Chýba už len samotné rozhodovacie pravidlo:
\[
    \text{pozorovanie} \rightarrow \text{politika} \rightarrow \text{akcia} \rightarrow \text{nové pozorovanie}.
\]
Agent teda vie vnímať aj konať, ale ešte nevieme, aká politika správania ho má riadiť.

Politiku správania budeme chápať ako funkciu
\[
    \pi_{\theta}(o_t) = a_t,
\]
kde \(o_t\) je pozorovanie, \(a_t\) akcia a \(\theta\) parametre politiky. V tejto kapitole nás už nezaujíma všeobecná definícia politiky, ale jej nájdenie: aké parametre majú viesť k jazde, ktorá prejde trať rýchlo a stabilne?

Ako reprezentáciu politiky použijeme viacvrstvovú neurónovú sieť. Tento výber nadväzuje na teoretickú časť: neurónová sieť je parametrická funkcia vhodná na aproximáciu nelineárnych vzťahov \cite{goodfellow2016deep} a dopredné siete s vhodnou nelinearitou majú silnú aproximačnú schopnosť \cite{cybenko1989approximation}. Pre nás je podstatné najmä to, že celá politika je potom určená konečným vektorom parametrov \(\theta\).

Tým sa problém nestráca, iba sa spresní. Namiesto ručne písaných pravidiel hľadáme bod v priestore parametrov siete. Každý bod predstavuje inú rozhodovaciu funkciu a v uzavretej slučke aj inú trajektóriu auta.

Postup kapitoly bude experimentálny. Najprv overíme, či sieť dokáže z navrhnutého pozorovania predikovať ľudskú akciu a aká architektúra na to stačí. Potom ukážeme limity samotného napodobňovania, navrhneme hodnotenie správania v prostredí a na ňom budeme stavať genetický tréning politiky. Vetvy učenia posilňovaním, Pareto optimalizácie a tréningových vylepšení zaradíme len tam, kde odpovedajú na konkrétny problém v tomto reťazci.

Prvé experimenty robíme v zjednodušenej rovinnej verzii \Trackmania{} bez výšky a povrchových zmien. Tak oddelíme samotné hľadanie politiky od ďalších komplikácií prostredia. Úvodná praktická otázka preto znie: akú architektúru neurónovej siete má zmysel ďalej optimalizovať?

\section{Výber architektúry neurónovej siete}

Ak má politika tvar \(\pi_{\theta}(o_t)=a_t\), potom architektúra siete určuje, koľko parametrov bude mať vektor \(\theta\). Teória nám hovorí, že dopredné siete s nelineárnymi aktiváciami sú silné aproximátory funkcií \cite{cybenko1989approximation,goodfellow2016deep}. V tejto časti však neriešime všeobecnú aproximačnú vetu. Potrebujeme prakticky zistiť, či na naše pozorovania stačí malá sieť, alebo už supervised dáta ukazujú potrebu väčšej kapacity.

Na tento účel sme použili záznamy jazdy hráča v zjednodušenej rovinnej verzii \Trackmania{} na asfaltovej mape. Každý vzor tvorí dvojica
\[
    (o_i, a_i^{\mathrm{human}}),
\]
kde \(o_i\) je pozorovanie agenta a \(a_i^{\mathrm{human}}\) akcia vykonaná hráčom. Tréning preto odpovedá na úzku otázku: vie daná sieť predikovať ľudskú akciu na zaznamenaných dátach? Ešte neodpovedá na otázku, či táto sieť dokáže samostatne dokončiť trať.

\begin{table}[htbp]
    \centering
    \small
    \renewcommand{\arraystretch}{1.14}
    \caption[Nastavenie sweep-u architektúr]{Nastavenie experimentu učenia s učiteľom použitého na porovnanie architektúr. Rovnaké dáta, stratová funkcia a tréningový rozpočet umožňujú porovnať siete podľa validačnej chyby a počtu parametrov.}
    \label{tab:supervised_architecture_setup}
    \begin{tabular}{>{\raggedright\arraybackslash}p{0.34\textwidth}
                    >{\raggedright\arraybackslash}p{0.56\textwidth}}
        \toprule
        Položka & Hodnota \\
        \midrule
        Tréningové vzorky & \(43\,898\) po filtrovaní a zrkadlení \\
        Validačné vzorky & \(10\,858\), bez zrkadlenia \\
        Vstup a výstup & \(34\) príznakov pozorovania, \(3\) akčné hodnoty \\
        Optimalizátor & Adam \\
        Stratová funkcia & \texttt{SmoothL1Loss}, váhy akcií \([1,1,3]\) \\
        Tréningový rozpočet & \(120\) epoch, batch size \(4096\) \\
        Učenie & learning rate \(0{,}001\), weight decay \(10^{-5}\) \\
        \bottomrule
    \end{tabular}
\end{table}

Do porovnania sme zaradili jednovrstvové aj dvojvrstvové siete, malé architektúry vhodné pre lacné experimenty a väčšie siete ako kapacitnú referenciu. Skúšali sme kombinácie aktivácií ReLU, \(\tanh\) a sigmoid. Sigmoid v skrytých vrstvách tu slúži najmä ako kontrolná vetva; pri tomto datasete dával výrazne horšie validačné straty než ReLU alebo kombinácie ReLU a \(\tanh\).

\begin{figure}[htbp]
    \centering
    \includegraphics[width=\textwidth]{training_policy/supervised_architecture_training_curves_all.pdf}
    \caption[Priebeh učenia všetkých architektúr]{Priebeh validačnej straty počas učenia všetkých porovnávaných architektúr. Ľavý panel ukazuje celý tréning, pravý panel približuje posledné epochy. Graf slúži na posúdenie stability učenia, nie na hodnotenie samostatnej jazdy agenta.}
    \label{fig:supervised_architecture_training_curves_all}
\end{figure}

Obrázok~\ref{fig:supervised_architecture_training_curves_all} ukazuje, že väčšina sietí sa neučila náhodne alebo nestabilne. Strata prudko klesla v prvých epochách a potom sa zlepšovanie spomaľovalo. Priblíženie konca tréningu zároveň ukazuje, že silnejšie modely sa ešte mierne zlepšovali aj ku koncu fixného rozpočtu. Výsledok preto čítame ako porovnanie pri rovnakom tréningovom rozpočte, nie ako dôkaz, že každá sieť už dosiahla svoje teoretické optimum.

\begin{figure}[htbp]
    \centering
    \includegraphics[width=0.92\textwidth]{training_policy/supervised_architecture_training_curves_practical.pdf}
    \caption[Priebeh učenia praktických kandidátov]{Detail praktických kandidátov po úvodnom poklese straty. Pri menších sieťach má ReLU,ReLU mierne nižšiu validačnú stratu než ReLU,\(\tanh\), no kapacitná referencia \(128\times64\) dosahuje najlepší výsledok práve s kombináciou ReLU,\(\tanh\).}
    \label{fig:supervised_architecture_training_curves_practical}
\end{figure}

Detail na obrázku~\ref{fig:supervised_architecture_training_curves_practical} ukazuje, že samotné dáta nerozhodujú jednoducho v prospech jednej aktivačnej kombinácie. Pri \(32\times16\) a \(48\times24\) vyšla o niečo lepšie kombinácia ReLU,ReLU. Pri väčšej referenčnej sieti \(128\times64\) však lepší výsledok dosiahla kombinácia ReLU,\(\tanh\). Preto by bolo prehnané tvrdiť, že supervised sweep sám dokazuje optimálnosť ReLU,\(\tanh\). Berieme ju ako pracovnú voľbu: kombinuje dve odlišné nelinearity a pri väčšej referencii nebola horšia, pričom skutočné rozhodnutie musí neskôr potvrdiť správanie v uzavretej slučke.

\begin{figure}[htbp]
    \centering
    \includegraphics[width=\textwidth]{training_policy/supervised_architecture_val_loss_vs_params.pdf}
    \caption[Sweep architektúr pri učení s učiteľom]{Porovnanie architektúr podľa najlepšej validačnej straty a počtu parametrov. Nižšia strata znamená lepšiu predikciu ľudskej akcie na validačných dátach, nie lepšiu samostatnú jazdu. Zvýraznené sú praktické kandidátske siete \(32\times16\) a \(48\times24\) s aktiváciami ReLU,\(\tanh\), väčšia sieť \(128\times64\) slúži ako kapacitná referencia.}
    \label{fig:supervised_architecture_val_loss_vs_params}
\end{figure}

Obrázok~\ref{fig:supervised_architecture_val_loss_vs_params} ukazuje očakávaný tradeoff: väčšie siete majú nižšiu validačnú stratu, ale rastie aj počet parametrov. Najnižšiu stratu dosiahla architektúra \(128\times64\) s aktiváciami ReLU,\(\tanh\), ale tá má \(12\,931\) parametrov. Pri neskoršej neuroevolúcii je každý parameter súčasťou prehľadávaného genómu, preto väčšia sieť neznamená automaticky lepšiu voľbu pre tréning agenta.

Pre praktické experimenty nás preto viac zaujímal kompromis okolo menších sietí. Architektúra \(32\times16\) má \(1\,699\) parametrov a ostáva lacným baseline modelom. Architektúra \(48\times24\) má \(2\,931\) parametrov a prináša lepšiu validačnú stratu bez toho, aby sa veľkosť genómu zväčšila rádovo.

\begin{table}[htbp]
    \centering
    \small
    \renewcommand{\arraystretch}{1.12}
    \caption[Vybrané architektúry zo sweep-u]{Vybrané architektúry z porovnania učenia s učiteľom. Tabuľka ukazuje, že väčší model znižuje validačnú stratu, ale zároveň zväčšuje priestor parametrov.}
    \label{tab:supervised_architecture_candidates}
    \begin{tabular}{llrrr}
        \toprule
        Architektúra & Aktivácie & Parametre & Val. strata & Steer MAE \\
        \midrule
        \(32\times16\) & ReLU, ReLU & \(1\,699\) & \(0{,}0659\) & \(0{,}2218\) \\
        \(32\times16\) & ReLU, \(\tanh\) & \(1\,699\) & \(0{,}0669\) & \(0{,}2247\) \\
        \(48\times24\) & ReLU, ReLU & \(2\,931\) & \(0{,}0606\) & \(0{,}2161\) \\
        \(48\times24\) & ReLU, \(\tanh\) & \(2\,931\) & \(0{,}0618\) & \(0{,}2163\) \\
        \(128\times64\) & ReLU, ReLU & \(12\,931\) & \(0{,}0578\) & \(0{,}2012\) \\
        \(128\times64\) & ReLU, \(\tanh\) & \(12\,931\) & \(0{,}0563\) & \(0{,}1993\) \\
        \bottomrule
    \end{tabular}
\end{table}

Z tohto experimentu si berieme dve pracovné architektúry. \(32\times16\) ReLU,\(\tanh\) používame ako lacný model na široké experimentovanie, kde je dôležitá rýchlosť porovnania metód. \(48\times24\) ReLU,\(\tanh\) používame ako silnejšieho kandidáta pre dlhšie behy. Porovnanie tým vyriešilo iba prvú časť problému: ukázalo, že navrhnuté pozorovanie a malá MLP sieť vedia napodobniť ľudskú akciu na dátach. Či sa takáto politika udrží aj pri vlastnej jazde, musí ukázať až správanie v prostredí.

\FloatBarrier

\section{Behaviorálne klonovanie v uzavretej slučke}

Ďalší test preto neprebiehal na validačnej množine, ale priamo na trati. Vybranú supervised politiku sme spustili v zjednodušenej rovinnej verzii \Trackmania{} na mape \texttt{single\_surface\_flat} a porovnali jej prejazd s desiatimi zaznamenanými jazdami učiteľa.

\begin{figure}[htbp]
    \centering
    \includegraphics[width=\textwidth]{training_policy/supervised_behavior_cloning_teacher_agent_paths.pdf}
    \caption[Trajektória supervised agenta oproti učiteľovi]{Porovnanie zaznamenaných jázd učiteľa a deterministického prejazdu politiky naučenej behaviorálnym klonovaním. Farba trajektórie vyjadruje rýchlosť v \(\mathrm{km/h}\). Agent trať dokončí, ale jeho trajektória je menej plynulá, odchyľuje sa od učiteľských prejazdov a obsahuje viac kontaktov so stenou.}
    \label{fig:supervised_behavior_cloning_teacher_agent_paths}
\end{figure}

Obrázok~\ref{fig:supervised_behavior_cloning_teacher_agent_paths} je dôležitý práve preto, že výsledok nie je ani úplné zlyhanie, ani hotové riešenie. Agent sa už z naučených dát dokáže dostať do cieľa. Zároveň však jazdí výrazne horšie než učiteľské prejazdy: na tejto mape dokončil jazdu za \(43{,}39\,\mathrm{s}\) a počas nej nazbieral \(32\) kontaktov so stenou, zatiaľ čo učiteľské jazdy v použitom datasete boli približne v rozsahu \(23{,}81\) až \(24{,}65\,\mathrm{s}\). Supervised strata teda zachytila schopnosť predikovať ľudskú akciu, ale sama o sebe nezaručila dobré správanie v uzavretej slučke.

Toto je presne miesto, kde sa behaviorálne klonovanie láme. Pri učení dostáva sieť pozorovania, ktoré vznikli počas ľudskej jazdy. Väčšina dobrých ľudských prejazdov však prirodzene vedie podobnou časťou trate: hráč sa snaží držať rýchlu stopu a vyhýba sa stavom, v ktorých by už musel prudko zachraňovať chybu. Dataset preto obsahuje veľa príkladov bežnej jazdy, ale málo príkladov zriedkavých opravných situácií.

Pri samostatnej jazde si agent svoje budúce pozorovania vytvára vlastnými akciami. Malá chyba v jednej zákrute posunie auto na inú stopu, tým zmení vzdialenosti k okrajom trate a ďalšie rozhodnutia už prebiehajú v stave, ktorý nemusel byť dobre pokrytý trénovacími dátami. Tento problém distribučného posunu je známy aj v literatúre k napodobňovaniu a motivuje interaktívne metódy typu DAgger \cite{ross2011dagger}.

Citlivosť na zásah do akcie sme si overili aj krátkou diagnostickou perturbáciou: na niekoľko krokov sme politike prepísali riadenie a sledovali sme, ako sa zmení zvyšok trajektórie. Tento beh nepoužívame ako samostatný kvantitatívny dôkaz, ale dobre ilustruje praktický problém: v sekvenčnej jazde sa odchýlka neuzavrie v jednom kroku. Mení budúcu polohu auta, a tým aj celé ďalšie rozhodovanie politiky.

\begin{figure}[htbp]
    \centering
    \includegraphics[width=\textwidth]{training_policy/supervised_behavior_cloning_butterfly_p04_left_120.pdf}
    \caption[Diagnostická zmena akcie pri supervised politike]{Diagnostická zmena akcie pri politike naučenej behaviorálnym klonovaním. Modrá čiara ukazuje pôvodný prejazd, oranžová prejazd po krátkom nanútenom zatočení doľava. Červený úsek označuje trvanie zásahu a čierne body kontakty so stenou. Zásah v tomto behu neviedol k horšiemu času, ale zmenil ďalšiu trajektóriu auta, čo ilustruje citlivosť jazdy na malé odchýlky v minulosti.}
    \label{fig:supervised_behavior_cloning_butterfly}
\end{figure}

Obrázok~\ref{fig:supervised_behavior_cloning_butterfly} používame ako analógiu k známej myšlienke citlivosti dynamických systémov na malé zmeny počiatočných podmienok, ktorá sa často označuje ako butterfly efekt \cite{lorenz1963deterministic}. Netvrdíme, že týmto diagnostickým behom dokazujeme chaotickosť nášho prostredia. Podstatná je jednoduchšia praktická pointa: v uzavretej slučke môže aj malá odchýlka dávno pred aktuálnym rozhodnutím dostať auto do inej časti stavového priestoru. Politika naučená iba na učiteľských dátach tam nemusí mať dostatok príkladov.

Z behaviorálneho klonovania si teda berieme pozitívny aj negatívny výsledok. Pozitívny výsledok je, že navrhnuté pozorovanie obsahuje dosť informácií na to, aby sieť vytvorila použiteľné akcie a prešla trať. Negatívny výsledok je, že samotné napodobňovanie nestačí na stabilného pretekárskeho agenta. Aby sme zachytili špecifické a opravné situácie, potrebovali by sme veľa variabilných vzoriek, alebo musíme dáta zbierať aj zo stavov, do ktorých sa agent dostane vlastnými chybami. Druhá možnosť vedie k interaktívnemu učeniu napodobňovaním.

\FloatBarrier

\section{Interaktívne učenie napodobňovaním}

Prirodzená oprava behaviorálneho klonovania je nezbierať iba ideálne jazdy človeka. Ak problém vzniká tým, že agent počas vlastnej jazdy navštívi iné stavy než učiteľ, potom má zmysel zbierať dáta práve v týchto stavoch. Tento smer nadväzuje na myšlienku DAgger: dataset sa postupne rozširuje o pozorovania, do ktorých sa dostáva aktuálna politika, zatiaľ čo expert poskytuje správnu akciu \cite{ross2011dagger}.

V našom interaktívnom režime preto človek a agent nejazdia oddelene. V každom kroku máme ľudskú akciu \(a_t^{\mathrm{human}}\) a akciu aktuálnej politiky \(a_t^{\mathrm{agent}}\). Do hry pošleme vykonanú akciu \(a_t^{\mathrm{exec}}\), ale ako trénovací cieľ ukladáme ľudskú korekciu. Pri režime miešania akcií môžeme vykonanú akciu zapísať ako
\[
    a_t^{\mathrm{exec}}
    =
    (1-\lambda_t)a_t^{\mathrm{human}}
    +
    \lambda_t a_t^{\mathrm{agent}},
\]
kde \(\lambda_t \in [0,1]\) určuje, akú časť riadenia v danom pokuse preberá agent. Na začiatku je \(\lambda_t\) malé a jazdu drží najmä človek. Postupne sa zvyšuje, takže agent vytvára čoraz väčšiu časť stavov, na ktorých sa bude po prejazde ďalej učiť.

Praktický cyklus jedného pokusu je krátky:
\[
\begin{array}{l}
\text{spusti aktuálnu politiku a nastav podiel agenta,}\\
\text{počas jazdy miešaj akciu človeka a agenta,}\\
\text{ukladaj pozorovanie a ľudskú korekčnú akciu,}\\
\text{po dokončení prejazdu krátko douč sieť na nových dátach,}\\
\text{zvýš podiel agenta a opakuj.}
\end{array}
\]
Tým sa od behaviorálneho klonovania líši najmä zdroj dát. Už nejde iba o stavy, do ktorých sa dostal dobrý ľudský prejazd. Do datasetu pribúdajú aj stavy vzniknuté tým, že agent zatáčal trochu inak, spomalil na inom mieste alebo vyšiel zo stopy. Práve tam sú korekčné akcie najcennejšie.

% TODO: doplniť graf priebehu interaktívneho učenia napodobňovaním: na osi x pokus alebo čas, na osi y podiel riadenia agenta \(\lambda\), doplnený o krátke doučenia po prejazdoch.

% TODO: doplniť obrázok trajektórií po interaktívnom učení napodobňovaním: aktuálny agent a predchádzajúce prejazdy/korekcie na rovnakej mape.

Tento krok však stále nemení cieľ učenia. Sieť sa učí napodobniť ľudskú korekciu v širšej množine stavov, nie priamo minimalizovať čas alebo počet kontaktov. Interaktívne napodobňovanie preto môže zlepšiť správanie oproti čisto statickému behaviorálnemu klonovaniu, ale ostáva závislé od kvality človeka, ktorý korekcie poskytuje. Ak hráč reaguje pomaly, volí slabú stopu alebo nevie ukázať dostatočne pestrú množinu opravných manévrov, agent sa bude učiť práve z týchto obmedzení.

Praktický problém je aj množstvo dát. Čím viac stavov chceme pokryť, tým viac prejazdov a korekcií potrebujeme. Pri krátkej mape sa to dá zvládnuť ručne, ale pri všeobecnejšom agentovi by sa počet potrebných vzoriek rýchlo zväčšoval. Po tejto časti teda vieme, že sieť vie konať a že vieme získavať lepšie korekčné dáta. Stále však potrebujeme spôsob, ako politiku hodnotiť podľa toho, čo na trati skutočne dosiahla.

\FloatBarrier
